%% toc: P15 161 161.2,161.4,161.5,161.9
% Proofs for the chapter on reading normal approximations and for the non-strict symmetric maximal inequality.
\section{Reading normal approximations: proofs}
\label{sec:annex-p15-w02c-p15-015-reading}

Use the notation of Chapter~\ref{sec:p15-w02c-p15-015-reading} and of
\cref{def:p15-clt-notation}. Weak convergence is convergence of CDFs at
every continuity point of the limiting CDF
(\cref{def:p15-weak-law-convergence}).

\begin{annexitem}{annex:p15-w02c-p15-015-chi-square}{thm:p15-w02c-p15-015-chi-square}{The density of the square}
\begin{proof}
Let \(F(y)=\mathbb P(\chi^2\leq y)\). For \(y<0\), \(F(y)=0\). Since
\(\chi\) has a density, \(\mathbb P(\chi=0)=0\), so \(F(0)=0\). For
\(y>0\),
\[
 F(y)=\mathbb P(-\sqrt y\leq\chi\leq\sqrt y)
     =2\int_0^{\sqrt y}\phi_0(u)\,du,
\]
using the evenness of \(\phi_0\). The integrand is continuous, so by the
fundamental theorem of calculus and the chain rule \(F\) is
differentiable on \((0,\infty)\) with
\[
 F'(y)=2\phi_0(\sqrt y)\cdot\frac1{2\sqrt y}
      =(2\pi y)^{-1/2}e^{-y/2}=f(y).
\]
For \(0<\varepsilon<y\), \(F(y)-F(\varepsilon)=\int_\varepsilon^yf\).
As \(\varepsilon\downarrow0\), \(F(\varepsilon)\to\mathbb P(\chi^2\leq0)
=F(0)=0\) by continuity of probability from above, and
\(\int_\varepsilon^yf\uparrow\int_0^yf\) by monotone convergence. Hence
\(F(y)=\int_0^yf\) for every \(y>0\), and \(f\) is a density of
\(\chi^2\).

The chi-squared density with one degree of freedom is
\((2^{1/2}\Gamma(\tfrac12))^{-1}y^{-1/2}e^{-y/2}\) on \(y>0\). The
substitution \(y=u^2/2\) gives
\[
 \Gamma(\tfrac12)=\int_0^\infty y^{-1/2}e^{-y}\,dy
   =\sqrt2\int_0^\infty e^{-u^2/2}\,du
   =\sqrt2\cdot\tfrac12\sqrt{2\pi}=\sqrt\pi,
\]
since \(\int_{\mathbb R}\phi_0=1\) and \(\phi_0\) is even. Thus
\((2^{1/2}\Gamma(\tfrac12))^{-1}=(2\pi)^{-1/2}\), and the two
densities coincide.

For the mean, \(\frac{d}{du}(-\phi_0(u))=u\phi_0(u)\), so integration by
parts on \([-A,A]\) gives
\[
 \int_{-A}^Au^2\phi_0(u)\,du=\bigl[-u\phi_0(u)\bigr]_{-A}^{A}
   +\int_{-A}^A\phi_0(u)\,du
   =-2A\phi_0(A)+\int_{-A}^A\phi_0 .
\]
As \(A\to\infty\), \(A\phi_0(A)\to0\) and the last integral tends to
\(1\); the left side increases to \(\mathbb E\chi^2\) by monotone
convergence. Hence \(\mathbb E\chi^2=1\). Finally
\((\sigma\chi)^2=\sigma^2\chi^2\) pointwise.
\end{proof}
\end{annexitem}

\begin{annexitem}{annex:p15-w02c-p15-015-histogram}{thm:p15-w02c-p15-015-histogram}{Telescoping intervals and the missing strip}
\begin{proof}
Since \(S\) takes only integer values, \(\{S\leq k\}\) and
\(\{S\geq k+1\}\) partition the sample space, so their probabilities
sum to \(1\). By \cref{def:p15-w02c-p15-015-histogram-rule}, the
corrected approximation of \(\mathbb P(S\leq k)\) is \(G(k+\frac12)\),
and that of \(\mathbb P(S\geq k+1)\) is
\(1-G((k+1)-\frac12)=1-G(k+\frac12)\). Their sum is exactly \(1\).

The corrected approximation of \(\mathbb P(S=m)\) is
\(G(m+\frac12)-G(m-\frac12)\). For a positive integer \(N\), the sum
over \(-N\leq m\leq N\) telescopes to \(G(N+\frac12)-G(-N-\frac12)\),
which tends to \(1-0=1\) as \(N\to\infty\) because \(G\) is a CDF.

Without the correction, \(\mathbb P(S\leq k)\) is read as \(G(k)\) and
\(\mathbb P(S\geq k+1)=1-\mathbb P(S<k+1)\) as \(1-G(k+1)\), the
continuity of \(G\) making the distinction between \(<\) and \(\leq\)
at \(k+1\) immaterial for \(G\). Their sum is
\(1-(G(k+1)-G(k))\). For \(G(x)=\Phi((x-a)/b)\) with \(b>0\),
\[
 G(k+1)-G(k)=\int_{(k-a)/b}^{(k+1-a)/b}\phi_0(u)\,du>0,
\]
because \(\phi_0>0\) everywhere and the interval has positive length
\(1/b\). Hence the uncorrected sum is strictly less than \(1\) for every
integer \(k\).
\end{proof}
\end{annexitem}

\begin{annexitem}{annex:p15-w02c-p15-015-pairwise-contrast}{thm:p15-w02c-p15-015-pairwise-contrast}{Etemadi's law and the zero atom}
\begin{proof}
The \(X_j\) are identically distributed fair signs, so
\(\mathbb E|X_1|=1<\infty\) and \(\mathbb EX_1=0\). They are pairwise
independent by \cref{thm:p15-clt-pairwise-boundary}, so
\cref{thm:p15-etemadi-strong-law} gives \(S_n/n\to0\) almost surely.
For \(i\ne j\), pairwise independence and boundedness give
\(\mathbb EX_iX_j=\mathbb EX_i\,\mathbb EX_j=0\); with
\(\mathbb EX_j^2=1\),
\(\operatorname{Var}(S_n)=\mathbb ES_n^2=\sum_j\mathbb EX_j^2
+\sum_{i\ne j}\mathbb EX_iX_j=n\).

By \cref{thm:p15-clt-pairwise-boundary},
\(\mathbb P(S_{2^r}=0)=1-2^{-r}\), so
\[
 \mathbb P\bigl(-\tfrac12<S_{2^r}/2^{r/2}\leq\tfrac12\bigr)
   \geq\mathbb P(S_{2^r}=0)=1-2^{-r}\longrightarrow1 .
\]
Suppose \(S_n/\sqrt n\Rightarrow m+s\chi\) with \(m\in\mathbb R\),
\(s>0\). The limiting CDF \(x\mapsto\Phi((x-m)/s)\) is continuous, so
convergence holds at \(\pm\frac12\), along every subsequence, in
particular along \(n=2^r\). Hence
\(\mathbb P(-\frac12<S_{2^r}/2^{r/2}\leq\frac12)\) would converge to
\(\Phi((\frac12-m)/s)-\Phi((-\frac12-m)/s)\), which is strictly less than
\(1\) because the normal density is positive on \((\frac12,\infty)\).
This contradicts the display. The case \(m=0\), \(s=1\) is the law of
\(\chi\), for which
\[
 \mathbb P(|\chi|\leq\tfrac12)=\int_{-1/2}^{1/2}\phi_0
   \leq\sup\phi_0=(2\pi)^{-1/2}<0.4,
\]
the last inequality because \(2\pi>6.25=0.4^{-2}\).
\end{proof}
\end{annexitem}

\begin{annexitem}{annex:p15-w02c-p15-015-truncation-window}{thm:p15-w02c-p15-015-truncation-window}{The tail, the truncated moment, and the empty Lindeberg events}
\begin{proof}
The hypothesis gives \(\mathbb P(|X_1|>1)=1\), and \(|X_1|\) has CDF
\(1-x^{-2}\) on \([1,\infty)\), hence density \(2x^{-3}\) there.

(1) \(\{Y_{n,m}\ne X_m\}=\{|X_m|>c_n\}\). By the union bound and
identical distribution,
\(\mathbb P(\bigcup_{m\leq n}\{|X_m|>c_n\})\leq n\mathbb P(|X_1|>c_n)
=n/c_n^2\), using \(c_n\geq1\). The quantity \(n/c_n^2=(c_n/\sqrt n)^{-2}\)
tends to zero exactly when \(c_n/\sqrt n\to\infty\).

(2) \(Y_{n,1}\) is bounded, and \(x\mapsto x\mathbf1_{\{|x|\leq c_n\}}\)
is odd, so the symmetry of \(X_1\) gives
\(\mathbb EY_{n,1}=-\mathbb EY_{n,1}=0\). Since \(Y_{n,1}^2\) is a
function of \(|X_1|\),
\[
 \mathbb EY_{n,1}^2=\int_1^{c_n}x^2\cdot2x^{-3}\,dx=2\log c_n,
\]
which increases with \(c_n\). Then
\(2\log c_n/\log n=1+2\log(c_n/\sqrt n)/\log n\), which tends to \(1\)
exactly when \(\log(c_n/\sqrt n)=o(\log n)\). For
\(c_n=n^{1/2+\varepsilon}\), \(2\log c_n=(1+2\varepsilon)\log n\).

(3) Let \(d_n=\sqrt{n\log n}\) (\(n\geq2\)) and \(\varepsilon>0\). The
Lindeberg sum of the row \(Y_{n,m}/d_n\) at level \(\varepsilon\) is
\(\sum_{m\leq n}\mathbb E(Y_{n,m}^2/d_n^2;|Y_{n,m}|>\varepsilon d_n)\).
Since \(|Y_{n,m}|\leq c_n\) and \(c_n/d_n\to0\), for all large \(n\) we
have \(c_n<\varepsilon d_n\), so every event in the sum is empty and the
sum is zero.

(4) For \(n\geq16\), \(\log n\geq\log16>e\), so \(\log\log n>1\) and
\(\log\log\log n>0\); also \(\sqrt n\geq4\), so \(c_n\geq1\). Then
\(n/c_n^2=(\log\log n)^{-2}\to0\);
\(\log(c_n/\sqrt n)=\log\log\log n\). Write \(u=\log n\geq e\), so
\(u\to\infty\). The inequality \(\log w\leq w^{1/2}\) for \(w\geq1\),
proved after this proof, gives \(0<\log\log u\leq\log u\leq u^{1/2}\),
hence \(\log\log\log n/\log n=\log\log u/u\leq u^{-1/2}\to0\). Applied
to \(w=\sqrt u\) it gives \(\log u=2\log\sqrt u\leq2u^{1/4}\), hence
\(c_n/\sqrt{n\log n}=\log u/\sqrt u\leq2u^{-1/4}\to0\). Items 1--3
apply.

(5) For \(c_n=n\): \(n/c_n^2=1/n\to0\), \(\mathbb EY_{n,1}^2=2\log n\)
by item~2, and \(c_n/\sqrt{n\log n}=\sqrt{n/\log n}\to\infty\), since
\(\log n\leq\sqrt n\) for \(n\geq1\) gives
\(n/\log n\geq\sqrt n\).
\end{proof}
The inequality \(\log w\leq w^{1/2}\) for \(w\geq1\) used twice above
holds because \(h(w)=w^{1/2}-\log w\) has \(h(1)=1>0\) and
\(h'(w)=\frac12w^{-1/2}-w^{-1}\), which is negative on \([1,4)\) and
nonnegative on \([4,\infty)\), so \(h\) attains its minimum on
\([1,\infty)\) at \(w=4\), where \(h(4)=2-\log4>0\) since \(4<e^2\).
\end{annexitem}

\begin{annexitem}{annex:p15-w02c-p15-015-power-tail-limit}{thm:p15-w02c-p15-015-power-tail-limit}{Scaling the characteristic function of a power tail}
\begin{proof}
\emph{The law of \(|X_1|\).} The hypothesis gives
\(\mathbb P(|X_1|>1)=1\), and \(|X_1|\) has CDF \(1-x^{-\alpha}\) on
\([1,\infty)\), hence density \(\alpha x^{-\alpha-1}\) there.

\emph{The characteristic function.} Let
\(\varphi(t)=\mathbb Ee^{itX_1}\). Since \(X_1\) and \(-X_1\) have the
same law, \(\varphi(t)=\varphi(-t)=\overline{\varphi(t)}\), so
\(\varphi\) is real and \(\varphi(t)=\mathbb E\cos(tX_1)
=\mathbb E\cos(t|X_1|)\), cosine being even. Hence, as
\(\int_1^\infty\alpha x^{-\alpha-1}dx=1\),
\[
 1-\varphi(t)=\int_1^\infty(1-\cos tx)\,\alpha x^{-\alpha-1}\,dx .
\]

\emph{The constant.} The function \(g(y)=y^2/2-1+\cos y\) has
\(g(0)=0\) and \(g'(y)=y-\sin y\geq0\) for \(y\geq0\), so
\(0\leq1-\cos y\leq y^2/2\); also \(1-\cos y\leq2\). Therefore
\[
 0\leq I_\alpha\leq\int_0^1\frac{y^2}2\alpha y^{-\alpha-1}\,dy
   +\int_1^\infty2\alpha y^{-\alpha-1}\,dy
   =\frac{\alpha}{2(2-\alpha)}+2<\infty,
\]
the first integral being finite because \(2-\alpha>0\). On
\([\pi/2,\pi]\), \(1-\cos y\geq1\), so
\(I_\alpha\geq\int_{\pi/2}^\pi\alpha y^{-\alpha-1}dy>0\). For real
\(t\ne0\), the substitution \(u=|t|y\) and the evenness of cosine give
\[
 \int_0^\infty(1-\cos ty)\,\alpha y^{-\alpha-1}\,dy=|t|^\alpha I_\alpha ;
\]
for \(t=0\) both sides vanish.

\emph{The scaled limit.} Fix \(t\) and put \(s_n=tn^{-1/\alpha}\). The
substitution \(x=n^{1/\alpha}y\), for which
\(\alpha x^{-\alpha-1}dx=n^{-1}\alpha y^{-\alpha-1}dy\) and
\(s_nx=ty\), gives
\[
 n\bigl(1-\varphi(s_n)\bigr)
   =\int_{n^{-1/\alpha}}^\infty(1-\cos ty)\,\alpha y^{-\alpha-1}\,dy .
\]
The lower limit decreases to \(0\) and the integrand is nonnegative,
so by monotone convergence the right side increases to
\(|t|^\alpha I_\alpha\).

\emph{The sum.} By independence and identical distribution,
\(\mathbb E\exp(itS_n/n^{1/\alpha})=\varphi(s_n)^n=(1+c_n/n)^n\) with
\(c_n=-n(1-\varphi(s_n))\to-I_\alpha|t|^\alpha\). By
\cref{thm:p15-clt-complex-products},
\(\mathbb E\exp(itS_n/n^{1/\alpha})\to\exp(-I_\alpha|t|^\alpha)\) for
every real \(t\). The limit is continuous at \(0\), so
\cref{thm:p15-characteristic-continuity-theorem} gives a probability law
\(Q_\alpha\) with this characteristic function and
\(S_n/n^{1/\alpha}\Rightarrow Q_\alpha\). No centering occurred in the
argument, and no moment of \(X_1\) was used.

\emph{Not normal.} By \cref{thm:p15-characteristic-normal-transform}, the
law \(N(m,s^2)\) of \(m+s\chi\), \(s\geq0\), has characteristic function
\(\exp(imt-s^2t^2/2)\), whose modulus is \(\exp(-s^2t^2/2)\). If
\(Q_\alpha\) were this law, the moduli of the two characteristic
functions would agree: \(I_\alpha|t|^\alpha=s^2t^2/2\) for all real
\(t\). Taking \(t=1\) gives \(s^2=2I_\alpha\); taking \(t=2\) then gives
\(2^\alpha I_\alpha=4I_\alpha\), so \(2^\alpha=4\) because
\(I_\alpha>0\), that is \(\alpha=2\), contrary to \(\alpha<2\). The case
\(s=0\) is included, and \(m=0\), \(s=1\) is the law of \(\chi\).
\end{proof}
\end{annexitem}

\begin{annexitem}{annex:p15-w02c-p15-015-rate-series}{thm:p15-w02c-p15-015-rate-series}{Comparison with a convergent power series}
\begin{proof}
Let \(0<\delta<1\) and \(p=(3-\delta)/2\), so \(p>1\). Then
\(n^{-1+\delta/2}\Delta_n\leq Cn^{-p}\). For \(n\geq2\),
\(n^{-p}\leq\int_{n-1}^nx^{-p}\,dx\) because \(x^{-p}\) is decreasing,
so \(\sum_{n\geq2}n^{-p}\leq\int_1^\infty x^{-p}\,dx=1/(p-1)<\infty\).
Hence \(\sum_nn^{-1+\delta/2}\Delta_n\leq C(1+1/(p-1))<\infty\).
For the second statement, if \(\sqrt n\,\Delta_n\leq C\) for all \(n\),
the first statement makes the series converge for every
\(\delta\in(0,1)\); so divergence for one such \(\delta\) forces
\(\sqrt n\,\Delta_n\) to be unbounded.
\end{proof}
\end{annexitem}

\begin{annexitem}{annex:p15-w02c-p15-015-levy-nonstrict}{thm:p15-w02c-p15-015-levy-nonstrict}{First crossing at the closed threshold}
\begin{proof}
Fix \(t>0\). For \(1\leq k\leq n\) define the first-crossing events
\[
 A_k=\{|T_k|\geq t,\ |T_j|<t\text{ for every }1\leq j<k\}.
\]
Each \(A_k\) is determined by \((Y_1,\ldots,Y_k)\). The \(A_k\) are
disjoint, because on \(A_k\) the index \(k\) is the first \(j\) with
\(|T_j|\geq t\); and their union is
\(\{\max_{1\leq k\leq n}|T_k|\geq t\}\), because on that event the first
such index exists and lies in \(\{1,\ldots,n\}\).

\emph{Splitting by sign.} Put \(A_k^+=A_k\cap\{T_k\geq t\}\) and
\(A_k^-=A_k\cap\{T_k\leq-t\}\). Since \(t>0\), these are disjoint, and
their union is \(A_k\) because \(|T_k|\geq t\) on \(A_k\).

\emph{The future sum.} Put \(R_k=T_n-T_k=\sum_{i=k+1}^nY_i\), with
\(R_n=0\). By \cref{thm:p15-independent-blocks}, \(R_k\) is independent
of the sigma-algebra generated by \(Y_1,\ldots,Y_k\), hence of
\(A_k^+\) and \(A_k^-\). The vector \((Y_{k+1},\ldots,Y_n)\) has the
product of its coordinate laws as its law, by independence; so does
\((-Y_{k+1},\ldots,-Y_n)\), and its coordinate laws are the same,
because each \(Y_i\) is symmetric. The two vectors therefore have the
same law, and so do their coordinate sums \(R_k\) and \(-R_k\). Hence
\(\mathbb P(R_k\geq0)=\mathbb P(-R_k\geq0)=\mathbb P(R_k\leq0)\). As
\(\{R_k\geq0\}\cup\{R_k\leq0\}\) is the whole space, these two equal
probabilities sum to at least \(1\), so each is at least \(\frac12\).
For \(k=n\) both equal \(1\).

\emph{Terminal inclusions.} On \(A_k^+\cap\{R_k\geq0\}\),
\(T_n=T_k+R_k\geq t\); on \(A_k^-\cap\{R_k\leq0\}\),
\(T_n=T_k+R_k\leq-t\). In both cases \(|T_n|\geq t\). The two events are
disjoint subsets of \(A_k\), so by independence
\[
 \begin{aligned}
 \mathbb P(A_k\cap\{|T_n|\geq t\})
 &\geq\mathbb P(A_k^+\cap\{R_k\geq0\})+\mathbb P(A_k^-\cap\{R_k\leq0\})\\
 &=\mathbb P(A_k^+)\,\mathbb P(R_k\geq0)+\mathbb P(A_k^-)\,\mathbb P(R_k\leq0)\\
 &\geq\tfrac12\bigl(\mathbb P(A_k^+)+\mathbb P(A_k^-)\bigr)
  =\tfrac12\mathbb P(A_k).
 \end{aligned}
\]

\emph{Summation.} The events \(A_k\cap\{|T_n|\geq t\}\), \(1\leq k\leq n\),
are disjoint subsets of \(\{|T_n|\geq t\}\). Summing over \(k\),
\[
 \mathbb P(|T_n|\geq t)\geq\sum_{k=1}^n\mathbb P(A_k\cap\{|T_n|\geq t\})
   \geq\tfrac12\sum_{k=1}^n\mathbb P(A_k)
   =\tfrac12\mathbb P\Bigl(\max_{1\leq k\leq n}|T_k|\geq t\Bigr),
\]
which is the claim. Only probabilities of events were used; no moment
of the summands is needed.

The two examples after the statement are direct evaluations. For
independent fair signs with \(n=2\) and \(t=1\), \(|T_1|=1\) always and
\(\mathbb P(|T_2|\geq1)=\mathbb P(T_2=\pm2)=\frac12\), so both sides equal
\(1\). For the constants \(Y_1=1\), \(Y_2=-1\), \(|T_1|=1\) and
\(T_2=0\), so the left side is \(1\) and the right side is \(0\).
\end{proof}
\end{annexitem}
